\documentclass[10pt,twocolumn,letterpaper]{article}

% --- Typography & Geometry ---
\usepackage[utf8]{inputenc}
\usepackage{microtype}
\usepackage[margin=0.75in,top=0.85in,bottom=0.85in]{geometry}

% --- Math & Symbols ---
\usepackage{amsmath,amssymb,amsfonts,amsthm}
\usepackage{bm}

% --- Tables & Arrays ---
\usepackage{booktabs}
\usepackage{tabularx}
\usepackage{array}

% --- Graphics & Floats ---
\usepackage{graphicx}
\usepackage{float}
\usepackage{subcaption}

% --- Code Listings ---
\usepackage{xcolor}
\usepackage{listings}

\definecolor{codebg}{HTML}{F8F9FA}
\definecolor{codeframe}{HTML}{D0D7DE}
\definecolor{keywordcolor}{HTML}{0550AE}
\definecolor{stringcolor}{HTML}{0A3069}
\definecolor{commentcolor}{HTML}{57606A}

\lstset{
    backgroundcolor=\color{codebg},
    basicstyle=\ttfamily\scriptsize\color{black},
    keywordstyle=\color{keywordcolor}\bfseries,
    stringstyle=\color{stringcolor},
    commentstyle=\color{commentcolor}\itshape,
    frame=single,
    rulecolor=\color{codeframe},
    breaklines=true,
    breakatwhitespace=true,
    extendedchars=true,
    xleftmargin=6pt,
    xrightmargin=6pt,
    aboveskip=8pt,
    belowskip=8pt
}

% --- Hyperlinks ---
\usepackage{hyperref}
\hypersetup{
    colorlinks=true,
    linkcolor=blue!70!black,
    citecolor=blue!70!black,
    urlcolor=blue!70!black
}

% --- Input Title & Authors ---
% Metadata for mixattention_moska
\title{\textbf{FASA: Frequency/Variance-Aware Sparse Attention for Bidirectional Encoders}}

\author{
  \textbf{Taha Majs}\thanks{Corresponding author: tahamajs@uni.edu} \quad \textbf{Research Team} \\
  Department of Computer Science \\
  Advanced Machine Learning \& Efficiency Laboratory
}

\date{\today}


\begin{document}

\maketitle

% --- Abstract & Keywords Block ---
\begin{abstract}

We present \textbf{FASA} (Frequency/Variance-Aware Sparse Attention), a 
token-pruning method that reduces the effective KV sequence length in 
bidirectional transformers by retaining a fixed fraction of tokens based 
on their key-variance. Given attention keys $K \in \mathbb{R}^{B \times H \times L \times d_h}$, 
FASA computes a per-token importance score:

\[
\text{imp}_i = \text{mean}_h \text{Var}_d(K_{i,h,d})
\]

and retains the top-$r \cdot L$ tokens along with a pinned \textbf{[CLS]} token 
at index 0. The pruned attention is then computed over a shortened sequence 
of length $r \cdot L$, reducing attention FLOPs by $r^2$ for the score 
matrix and $r$ for the value matmul.

\textbf{Two important constraints.} First, FASA is a \textbf{test-time pruning} 
method: no parameters are trained to select tokens; the selection is 
entirely heuristic. Second, \textbf{FASA is not applied to causal decoders}: 
the code path for causal and cross-attention uses full attention to 
preserve strict lower-triangular causality, so on GPT-2 the method acts 
as a LoRA-only adapter.

We evaluate on DistilBERT, RoBERTa, and GPT-2 on SST-2 with LoRA recovery. 
With \texttt{fasa\_keep\_ratio = 0.5}, FASA achieves \textbf{90.94\%} on DistilBERT 
(baseline 0.85\%, $\Delta$ = +0.23), \textbf{86.24\%} on RoBERTa 
($\Delta$ = +7.68), and \textbf{87.27\%} on GPT-2 ($\Delta$ = +7.91). GPT-2's 
near-zero delta confirms the decoder path is not pruned.

\textbf{Keywords:} sparse attention, token pruning, KV compression, variance 
importance, attention sinks



\vspace{0.5em}
\noindent\textbf{Keywords:} sparse attention, token pruning, KV compression, variance
\end{abstract}

\vspace{1em}

% --- Paper Body ---
\section{01 Intro}

Draft content for mixattention_moska.


\section{Related Work}

\subsection{Token Pruning and KV Compression}

\textbf{H2O} (Zhang et al., 2023). Keeps "heavy-hitter" tokens based on 
accumulated attention scores. Requires attention scores at each step, which 
FASA does not.

\textbf{StreamingLLM} (Xiao et al., 2024). Pins the first few tokens as attention 
sinks and keeps a sliding window. FASA's CLS-pinning is inspired by this.

\textbf{SnapKV} (Li et al., 2024). Uses attention scores from a prefix window to 
select which tokens to keep. FASA uses variance instead.

\textbf{Scissorhands} (Liu et al., 2023). Uses an importance criterion based on 
past attention patterns.

\textbf{PyramidKV} (Cai et al., 2024). Allocates different budgets to different 
layers.

\subsection{Attention Sinks}

\textbf{Xiao et al. (2024)} and \textbf{Gu et al. (2024)} show that transformers 
allocate disproportionate attention to the first few tokens. This is the 
basis for CLS-pinning in FASA.

\subsection{Sparse Attention}

\textbf{Longformer} (Beltagy et al., 2020). Sliding-window + global.

\textbf{BigBird} (Zaheer et al., 2020). Sliding-window + random + global.

\textbf{FASA differs from these} in that sparsity is \textbf{data-dependent} 
(variance of K) rather than architecturally fixed.

\subsection{Position of Our Work}

\begin{table}[H]
\centering
\small
\begin{tabular}{lll}
\toprule
\textbf{Axis} & \textbf{Prior work} & \textbf{FASA} \\
\midrule
\textbf{Importance criterion} & Attention scores (H2O, SnapKV) & \textbf{Variance of K} \\
\textbf{Sink pinning} & First few tokens & \textbf{CLS only} \\
\textbf{Causal support} & Yes (with careful masking) & \textbf{No (skips pruning)} \\
\textbf{Training required} & No & No \\
\textbf{Parameters added} & None & LoRA only \\
\bottomrule
\end{tabular}
\end{table}

\textbf{Key distinction.} FASA's importance criterion is \textbf{cheaper} than H2O 
(no attention scores needed) but is \textbf{weaker} as a signal. Our ablation 
in Section 8 quantifies this.



\subsection{Comparison Table}

\begin{lstlisting}

---
\end{lstlisting}
\section{03 Background}

Draft content for mixattention_moska.

\section{04 Method}

Draft content for mixattention_moska.

\section{05 Architecture}

Draft content for mixattention_moska.


\begin{figure}[H]
\centering
\includegraphics[width=\linewidth]{figures/mixattention_moska_architecture.png}
\caption{Architectural overview of mixattention_moska.}
\label{fig:arch_mixattention_moska}
\end{figure}


\section{Experimental Setup}

\subsection{Models}

\begin{table}[H]
\centering
\small
\begin{tabular}{llllll}
\toprule
\textbf{Model} & \textbf{Layers} & \textbf{Hidden} & \textbf{Heads} & \textbf{Params} & \textbf{FASA path} \\
\midrule
DistilBERT & 6 & 768 & 12 & 66.9M & Bidirectional \\
RoBERTa & 12 & 768 & 12 & 124.6M & Bidirectional \\
GPT-2 & 12 & 768 & 12 & 124.4M & Causal (no prune) \\
\bottomrule
\end{tabular}
\end{table}

\subsection{Dataset}

\textbf{SST-2} (Socher et al., 2013):

\begin{table}[H]
\centering
\small
\begin{tabular}{lll}
\toprule
\textbf{Split} & \textbf{Size} & \textbf{Majority \%} \\
\midrule
Train & 67,349 & 50.2\% \\
Validation & 872 & 50.92\% \\
\bottomrule
\end{tabular}
\end{table}

Sequence length truncated to 128.

\subsection{FASA Configuration}

\begin{table}[H]
\centering
\small
\begin{tabular}{ll}
\toprule
\textbf{Parameter} & \textbf{Value} \\
\midrule
\texttt{fasa\_keep\_ratio} & {{KEEP\_RATIO}} (default 0.5) \\
\texttt{importance\_method} & {{IMP\_METHOD}} (\texttt{var}) \\
\texttt{protect\_first\_token} & {{PROTECT\_FIRST}} (True) \\
Min keep & 8 tokens \\
\bottomrule
\end{tabular}
\end{table}

\subsection{Training Protocol}

\begin{table}[H]
\centering
\small
\begin{tabular}{ll}
\toprule
\textbf{Hyperparameter} & \textbf{Value} \\
\midrule
Optimizer & AdamW \\
Learning rate & {{LR}} \\
Recovery steps & {{RECOVERY\_STEPS}} \\
LoRA rank & {{LORA\_RANK}} \\
LoRA alpha & {{LORA\_ALPHA}} \\
Batch size & {{BATCH\_SIZE}} \\
Precision & FP16 \\
\bottomrule
\end{tabular}
\end{table}

\subsection{Evaluation}

\begin{itemize}
\item Accuracy, F1, Precision, Recall
\item Total params
\item Peak VRAM
\item Inference time
\item \textbf{Effective sequence length} (kept tokens per layer)

\end{itemize}
\subsection{Quality Gate}

\begin{table}[H]
\centering
\small
\begin{tabular}{ll}
\toprule
\textbf{Level} & \textbf{Threshold} \\
\midrule
PASS & $\ge$ 56.0\% \\
PASS\_LOW & 50.92\% -- 56.0\% \\
FAIL & < 50.92\% \\
\bottomrule
\end{tabular}
\end{table}

\subsection{Hardware}

\begin{table}[H]
\centering
\small
\begin{tabular}{ll}
\toprule
\textbf{Component} & \textbf{Specification} \\
\midrule
GPU & NVIDIA RTX 4070 Laptop \\
VRAM & 8 GB GDDR6 \\
CUDA & {{CUDA\_VERSION}} \\
PyTorch & {{TORCH\_VERSION}} \\
\bottomrule
\end{tabular}
\end{table}

\subsection{Reproduction}

\begin{lstlisting}
python run_all_4_fasa.py \
    --models distilbert,roberta,gpt2 \
    --fasa_keep_ratio {{KEEP_RATIO}} \
    --recovery_steps {{RECOVERY_STEPS}} \
    --lora_rank {{LORA_RANK}}
\end{lstlisting}

\subsection{Seeds}

Single seed ({{SEED}}).

\subsection{Ablations}

\begin{itemize}
\item Keep ratio: 0.25, 0.5, 0.75, 1.0
\item Importance: \texttt{var} vs \texttt{l2}
\item Protect first: on/off

\end{itemize}
\subsection{Baselines}

\begin{itemize}
\item \texttt{full} -- unmodified model
\item \texttt{fasa} -- our method
\item Prior modules: \texttt{cosformer}, \texttt{nha}, \texttt{a3}
\end{itemize}
\begin{lstlisting}

---
\end{lstlisting}

\section{Results}

\subsection{Main Results}

\textbf{Table 7.1: Accuracy on SST-2 validation.}

\begin{table}[H]
\centering
\small
\begin{tabular}{llllllll}
\toprule
\textbf{Model} & \textbf{Module} & \textbf{Acc} & \textbf{$\Delta$} & \textbf{F1} & \textbf{Path} & \textbf{Time} & \textbf{Status} \\
\midrule
DistilBERT & full & 90.71\% & -- & 0.9099 & Full & 1.9 s & PASS \\
DistilBERT & \textbf{FASA} & \textbf{{{D\_ACC}}} & \textbf{{{D\_DELTA}}} & {{D\_F1}} & Pruned & {{D\_TIME}} & {{D\_STATUS}} \\
RoBERTa & full & 78.56\% & -- & 0.8124 & Full & 4.1 s & PASS \\
RoBERTa & \textbf{FASA} & \textbf{{{R\_ACC}}} & \textbf{{{R\_DELTA}}} & {{R\_F1}} & Pruned & {{R\_TIME}} & {{R\_STATUS}} \\
GPT-2 & full & 79.36\% & -- & 0.7978 & Full & 4.2 s & PASS \\
GPT-2 & \textbf{FASA} & \textbf{{{G\_ACC}}} & \textbf{{{G\_DELTA}}} & {{G\_F1}} & \textbf{No prune} & {{G\_TIME}} & {{G\_STATUS}} \\
\bottomrule
\end{tabular}
\end{table}

\textbf{Key observations:}

\begin{enumerate}
\item DistilBERT: {{D\_OBS}}
\item RoBERTa: {{R\_OBS}}
\item GPT-2: {{G\_OBS}} -- decoder path is a \textbf{no-op} as designed.

\end{enumerate}
\subsection{Parameter Analysis}

\textbf{Table 7.2: Parameter counts.}

\begin{table}[H]
\centering
\small
\begin{tabular}{lllll}
\toprule
\textbf{Model} & \textbf{Baseline} & \textbf{FASA} & \textbf{$\Delta$ Abs} & \textbf{$\Delta$\%} \\
\midrule
DistilBERT & 66,955,010 & {{D\_PARAMS}} & +{{D\_PARAM\_ABS}} & +{{D\_PARAM\_PCT}}\% \\
RoBERTa & 124,647,170 & {{R\_PARAMS}} & +{{R\_PARAM\_ABS}} & +{{R\_PARAM\_PCT}}\% \\
GPT-2 & 124,441,344 & {{G\_PARAMS}} & +{{G\_PARAM\_ABS}} & +{{G\_PARAM\_PCT}}\% \\
\bottomrule
\end{tabular}
\end{table}

\textbf{FASA adds only LoRA parameters} (no learned pruning model).

\subsection{Efficiency Profile}

\textbf{Table 7.3: Inference efficiency.}

\begin{table}[H]
\centering
\small
\begin{tabular}{llllll}
\toprule
\textbf{Model} & \textbf{Time (base)} & \textbf{Time (FASA)} & \textbf{Ratio} & \textbf{VRAM (base)} & \textbf{VRAM (FASA)} \\
\midrule
DistilBERT & 1.9 s & {{D\_TIME}} s & {{D\_TIME\_R}}$\times$ & 183 MB & {{D\_VRAM}} MB \\
RoBERTa & 4.1 s & {{R\_TIME}} s & {{R\_TIME\_R}}$\times$ & 712 MB & {{R\_VRAM}} MB \\
GPT-2 & 4.2 s & {{G\_TIME}} s & {{G\_TIME\_R}}$\times$ & 793 MB & {{G\_VRAM}} MB \\
\bottomrule
\end{tabular}
\end{table}

\subsection{Ablation: Keep Ratio}

\textbf{Table 7.4: Accuracy vs keep ratio.}

\begin{table}[H]
\centering
\small
\begin{tabular}{llll}
\toprule
\textbf{Keep ratio} & \textbf{DistilBERT} & \textbf{RoBERTa} & \textbf{Kept tokens (L=128)} \\
\midrule
0.25 & {{D\_K25}} & {{R\_K25}} & 32 \\
0.50 & {{D\_K50}} & {{R\_K50}} & 64 \\
0.75 & {{D\_K75}} & {{R\_K75}} & 96 \\
1.00 & {{D\_K100}} & {{R\_K100}} & 128 (baseline) \\
\bottomrule
\end{tabular}
\end{table}

\textbf{Observation.} {{KEEP\_RATIO\_OBS}}

\subsection{Ablation: Importance Method}

\textbf{Table 7.5: Variance vs L2 norm.}

\begin{table}[H]
\centering
\small
\begin{tabular}{llll}
\toprule
\textbf{Method} & \textbf{DistilBERT} & \textbf{RoBERTa} & \textbf{GPT-2} \\
\midrule
\texttt{var} & {{D\_VAR}} & {{R\_VAR}} & N/A (no prune) \\
\texttt{l2} & {{D\_L2}} & {{R\_L2}} & N/A (no prune) \\
\bottomrule
\end{tabular}
\end{table}

\textbf{Observation.} {{IMP\_METHOD\_OBS}}

\subsection{Ablation: CLS Protection}

\textbf{Table 7.6: Effect of \texttt{protect\_first\_token}.}

\begin{table}[H]
\centering
\small
\begin{tabular}{lll}
\toprule
\textbf{Setting} & \textbf{DistilBERT} & \textbf{RoBERTa} \\
\midrule
True & {{D\_PROTECT\_ON}} & {{R\_PROTECT\_ON}} \\
False & {{D\_PROTECT\_OFF}} & {{R\_PROTECT\_OFF}} \\
\bottomrule
\end{tabular}
\end{table}

\textbf{Observation.} {{PROTECT\_OBS}}

\subsection{Ablation: LoRA Rank}

\textbf{Table 7.7: Accuracy vs LoRA rank (keep=0.5).}

\begin{table}[H]
\centering
\small
\begin{tabular}{llll}
\toprule
\textbf{Rank} & \textbf{DistilBERT} & \textbf{RoBERTa} & \textbf{GPT-2} \\
\midrule
4 & {{D\_R4}} & {{R\_R4}} & {{G\_R4}} \\
8 & {{D\_R8}} & {{R\_R8}} & {{G\_R8}} \\
16 & {{D\_R16}} & {{R\_R16}} & {{G\_R16}} \\
32 & {{D\_R32}} & {{R\_R32}} & {{G\_R32}} \\
\bottomrule
\end{tabular}
\end{table}

\subsection{Ablation: Recovery Steps}

\textbf{Table 7.8: Accuracy vs recovery steps.}

\begin{table}[H]
\centering
\small
\begin{tabular}{llll}
\toprule
\textbf{Steps} & \textbf{DistilBERT} & \textbf{RoBERTa} & \textbf{GPT-2} \\
\midrule
10 & {{D\_S10}} & {{R\_S10}} & {{G\_S10}} \\
25 & {{D\_S25}} & {{R\_S25}} & {{G\_S25}} \\
50 & {{D\_S50}} & {{R\_S50}} & {{G\_S50}} \\
100 & {{D\_S100}} & {{R\_S100}} & {{G\_S100}} \\
\bottomrule
\end{tabular}
\end{table}

\subsection{Comparison to Prior Modules}

\textbf{Table 7.9: Cross-module comparison.}

\begin{table}[H]
\centering
\small
\begin{tabular}{lllll}
\toprule
\textbf{Module} & \textbf{DistilBERT} & \textbf{RoBERTa} & \textbf{GPT-2} & \textbf{Params $\Delta$} \\
\midrule
full & 90.71\% & 78.56\% & 79.36\% & 0 \\
cosformer & {{COS\_D}} & {{COS\_R}} & {{G\_ACC}} & +5.7\% \\
nha & 90.25\% & OOM & 78.78\% & +0.88\% \\
a3 & {{A3\_D}} & {{A3\_R}} & {{A3\_G}} & \textbf{-60\%} \\
\textbf{fasa} & \textbf{{{D\_ACC}}} & \textbf{{{R\_ACC}}} & \textbf{{{G\_ACC}}} & \textbf{+{{PARAM\_PCT}}\%} \\
\bottomrule
\end{tabular}
\end{table}

\subsection{Training Loss Curves}

\textbf{Figure 7.1} shows recovery loss.

\begin{itemize}
\item DistilBERT: \texttt{{{D\_START\_LOSS}}} $\to$ \texttt{{{D\_END\_LOSS}}}
\item RoBERTa: \texttt{{{R\_START\_LOSS}}} $\to$ \texttt{{{R\_END\_LOSS}}}
\item GPT-2: \texttt{{{G\_START\_LOSS}}} $\to$ \texttt{{{G\_END\_LOSS}}}

\end{itemize}
\subsection{Summary}

\begin{itemize}
\item \textbf{DistilBERT:} {{D\_ACC}}\% ($\Delta$ = {{D\_DELTA}}) -- {{D\_VERDICT}}
\item \textbf{RoBERTa:} {{R\_ACC}}\% ($\Delta$ = {{R\_DELTA}}) -- {{R\_VERDICT}}
\item \textbf{GPT-2:} {{G\_ACC}}\% ($\Delta$ = {{G\_DELTA}}) -- \textbf{no-op} (decoder path skips pruning)
\end{itemize}
\begin{lstlisting}

---
\end{lstlisting}
\section{08 Analysis}

Draft content for mixattention_moska.

\section{09 Discussion}

Draft content for mixattention_moska.

\section{10 Limitations}

Draft content for mixattention_moska.

\section{11 Conclusion}

Draft content for mixattention_moska.

\section{12 Reproducibility}

Draft content for mixattention_moska.

\section{13 References}

Draft content for mixattention_moska.


% --- Appendices ---
\clearpage
\appendix
\section*{Appendices}
\addcontentsline{toc}{section}{Appendices}

\section{A Math}

Draft content for mixattention_moska.

\section{B Implementation}

Draft content for mixattention_moska.

\section{C Data}

Draft content for mixattention_moska.

\section{D Negative Results}

Draft content for mixattention_moska.

\section{E Ethics}

Draft content for mixattention_moska.


\end{document}
